\chapter{Ambient Score Tangents and the Hermite Spectrum}
\label{ch:gaussian-hermite}

At a fixed law, three different objects are called tangent directions: the
tangent space of a finite-dimensional family, the Hilbert space of centered
square-integrable scores, and the bounded action quotient on which the action
calculus of \Cref{thm:rg-bounded-action-calculus} is analytic. This chapter asks
how the three are nested at a Gaussian law, and what a parameter-independent
block channel does to scores once a block of independent copies has been formed.
The main result, \Cref{thm:rg-gaussian-hermite-spectrum}, computes the score
operator of the normalized sum of \(b\) independent standard Gaussian coordinates:
in the normalized Hermite basis it is diagonal with eigenvalues \(b^{1-k/2}\),
compact, and injective with dense but proper range. It rests on the exact
extensive Fisher budget of an arbitrary parameter-independent block channel
(\Cref{thm:rg-score-pushforward-defect}) and on the realization of every
centered square-integrable score by a two-sided path
(\Cref{lem:rg-dqm-realization}), and
\Cref{prop:ig-hermite-exponential-domain} decides which Hermite directions also
admit two-sided exponential action charts. The chapter is an explicit instance of
a distinction drawn in general in the main manuscript. Exact arrows create no
growth of a tangent direction: the reverse conditional expectation is a
contraction (\Cref{thm:rg-action-lp-contraction}), the bounded action sector is
nonexpansive \eqref{eq:rg-nonlinear-action-sup-contraction}, and the exact
interaction cocycle is bounded uniformly in depth
(\Cref{prop:rg-interaction-native-growth-bound}). The eigenvalue \(b^{1/2}>1\) of
the degree-one direction therefore comes entirely from the extensive replication
factor \(\sqrt b\), a declared normalization between two reference tangent spaces.

The statements of the chapter are linear statements about score tangents at one
law, and its spectral statements concern the scalar independent normalized sum.
The chapter does not assert a convergence statement for the nonlinear map that
the normalized sum induces on laws, or a statement about the nonlinear bounded
action map beyond its derivative. \status{NOT-CLAIMED}

\section{Quadratic-mean scores}
\label{sec:rg-score-lift}

\definitionheading{Quadratic-mean differentiability and score}{def:rg-dqm-score}
Let \((\mathsf W,\mathscr W,\nu)\) be a probability space and write
\begin{equation}
L^2_0(\nu)
=\left\{h\in L^2(\nu):\int h d\nu=0\right\}.
\label{eq:rg-centered-tangent}
\end{equation}
A family \((\nu_t)_{\lvert t\rvert<t_0}\) of probability laws with
\(\nu_t\ll\nu\) and \(\nu_0=\nu\) is differentiable in quadratic mean at \(t=0\)
with score \(h\in L^2(\nu)\) when, for \(t\to0\) from both sides,
\begin{equation}
\left\lVert
\sqrt{\frac{d\nu_t}{d\nu}}-1-\frac t2h
\right\rVert_{L^2(\nu)}
=o(\lvert t\rvert).
\label{eq:rg-dqm-definition}
\end{equation}
The Fisher information of the path at \(t=0\) is \(\lVert h\rVert_{L^2(\nu)}^2\).

\lemmaheading{Every centered square-integrable score is realized two-sided}{lem:rg-dqm-realization}
The score of every path that is differentiable in quadratic mean at \(\nu\) lies in
\(L^2_0(\nu)\). Conversely, let \(h\in L^2_0(\nu)\) and, for every real \(t\), put
\(Z_t=1+t^2\lVert h\rVert_2^2/4\) and
\begin{equation}
\frac{d\nu_t}{d\nu}=\frac{(1+th/2)^2}{Z_t}.
\label{eq:rg-quadratic-score-path}
\end{equation}
Then every \(\nu_t\) is a probability law, \(\nu_0=\nu\), and the path is
differentiable in quadratic mean at \(t=0\) with score \(h\), for \(t\to0^+\) and for
\(t\to0^-\). No exponential moment and no boundedness of \(h\) is used. The
quadratic-mean tangent set at \(\nu\) is therefore exactly \(L^2_0(\nu)\).
\status{ESTABLISHED}

\noindent\emph{Proof.}
For the first statement write \(p_t=d\nu_t/d\nu\). Since \(\int p_td\nu=1\),
\(\int(\sqrt{p_t}-1)d\nu=-\tfrac12\int(\sqrt{p_t}-1)^2d\nu\). By
\eqref{eq:rg-dqm-definition} the left side is \(\tfrac t2\int hd\nu+o(\lvert t\rvert)\)
and the right side is \(O(t^2)\), so \(\int hd\nu=0\).

For the converse, centering gives
\(\int(1+th/2)^2d\nu=1+t\nu(h)+\tfrac{t^2}4\nu(h^2)=Z_t\), so \(p_t\geq0\) integrates
to one. Put \(w_t=1+th/2\), so that \(\sqrt{p_t}=\lvert w_t\rvert/\sqrt{Z_t}\) and
\begin{equation}
\sqrt{p_t}-w_t
=\frac{\lvert w_t\rvert-w_t}{\sqrt{Z_t}}
+w_t\left(\frac1{\sqrt{Z_t}}-1\right).
\label{eq:rg-quadratic-path-split}
\end{equation}
On \(\{w_t<0\}\) one has \(th/2<-1\), hence \(\lvert h\rvert>2/\lvert t\rvert\) and
\(\lvert w_t\rvert-w_t=2\lvert w_t\rvert\leq\lvert th\rvert\); elsewhere the difference
vanishes. Since \(Z_t\geq1\),
\begin{equation}
\frac1{t^2}\left\lVert
\frac{\lvert w_t\rvert-w_t}{\sqrt{Z_t}}
\right\rVert_2^2
\leq\int h^2\mathbf 1_{\{\lvert h\rvert>2/\lvert t\rvert\}}d\nu
\longrightarrow0
\label{eq:rg-quadratic-path-tail}
\end{equation}
by dominated convergence, because \(h^2\) is integrable. For the second term of
\eqref{eq:rg-quadratic-path-split}, \(\lVert w_t\rVert_2=\sqrt{Z_t}\), so its norm is
exactly \(\lvert1-\sqrt{Z_t}\rvert=O(t^2)\). Adding the two bounds and using
\(w_t=1+\tfrac t2h\) gives \eqref{eq:rg-dqm-definition}. \hfill$\square$

\section{The extensive block lift and the pushed score}

\propositionheading{Independent block replication of a differentiable path}{prop:rg-score-block-lift}
Fix an integer \(b\geq1\). (i) Let \((\nu_t)\) satisfy \eqref{eq:rg-dqm-definition}
with score \(h\). Then the product path \(\nu_t^{\otimes b}\) is differentiable in
quadratic mean at \(\nu^{\otimes b}\) with score
\begin{equation}
(\mathscr I_bh)(x_1,\dots,x_b)=\sum_{i=1}^bh(x_i),
\label{eq:rg-score-replication}
\end{equation}
and the replication lift is a bounded linear injection
\begin{equation}
\mathscr I_b:L^2_0(\nu)\longrightarrow L^2_0(\nu^{\otimes b}),
\qquad
\lVert\mathscr I_bh\rVert_{L^2(\nu^{\otimes b})}^2
=b\lVert h\rVert_{L^2(\nu)}^2,
\label{eq:rg-score-replication-norm}
\end{equation}
so \(\lVert\mathscr I_b\rVert\leq\sqrt b\), with equality when \(L^2_0(\nu)\neq\{0\}\), and
\(b^{-1/2}\mathscr I_b\) is an isometry onto its range. The lift is not the Hoeffding assembly \(E_\ell\) of
\eqref{eq:rg-interaction-assembly-extraction}.
(ii) For \(b\geq2\), no parameter-independent Markov kernel \(\mathsf K\) satisfies
\(\mathcal N(x,1)\mathsf K=\mathcal N(x\mathbf 1_b,I_b)\) for every \(x\in\R\).
\status{ESTABLISHED}

\noindent\emph{Proof.}
(i) Write \(r_t=\sqrt{d\nu_t/d\nu}=1+\tfrac t2h+t\varepsilon_t\) with
\(\lVert\varepsilon_t\rVert_{L^2(\nu)}\to0\). The product amplitude is
\(\prod_{i=1}^b(1+\delta_i)\) with
\(\delta_i(x)=\tfrac t2h(x_i)+t\varepsilon_t(x_i)\), and
\begin{equation}
\prod_{i=1}^b(1+\delta_i)
=1+\sum_{i=1}^b\delta_i
+\sum_{\lvert S\rvert\geq2}\prod_{i\in S}\delta_i .
\label{eq:rg-product-amplitude-expansion}
\end{equation}
Each factor depends on one coordinate, so \(L^2(\nu^{\otimes b})\) norms of products
factor, \(\lVert\prod_{i\in S}\delta_i\rVert_2=\prod_{i\in S}\lVert\delta_i\rVert_2
=O(\lvert t\rvert^{\lvert S\rvert})\), and every term with \(\lvert S\rvert\geq2\) is
\(O(t^2)\). The linear part is \(\tfrac t2\mathscr I_bh+t\sum_i\varepsilon_t(x_i)\),
whose remainder has norm at most \(b\lvert t\rvert\lVert\varepsilon_t\rVert_2
=o(\lvert t\rvert)\). Independence and centering give
\(\E(\mathscr I_bh)^2=\sum_i\E h(x_i)^2=b\lVert h\rVert_2^2\).
(ii) The family \(x\mapsto\mathcal N(x,1)\) has Fisher information one and the family
\(x\mapsto\mathcal N(x\mathbf 1_b,I_b)\) has Fisher information \(b\). A
parameter-independent Markov kernel cannot increase Fisher information
(\Cref{thm:cg-fisher-contraction}), so such a \(\mathsf K\) would force \(b\leq1\).
\hfill$\square$

\medskip
The replication is extensive. Replacing the sum \eqref{eq:rg-score-replication} by the
arithmetic mean without rescaling the path parameter would change the Fisher
information from \(b\lVert h\rVert_2^2\) to \(\lVert h\rVert_2^2/b\), which is a different
construction. Part (ii) says that the replication pair is a lift between two
reference tangent spaces and not a pushforward along any arrow, so it can serve
neither as a counterexample to Markov contraction nor as evidence that Fisher
information increased along a coarse-graining arrow. The factor \(\sqrt b\) belongs to
the declared normalization that identifies the two reference spaces.

\theoremheading{Pushed score and the exact extensive Fisher defect}{thm:rg-score-pushforward-defect}
Let \(C_b:\mathsf W^b\rightsquigarrow\mathsf Z\) be a parameter-independent Markov
kernel between standard Borel spaces, put \(\nu^{(b)}=\nu^{\otimes b}\),
\(\mathbb J_b(dx,dz)=\nu^{(b)}(dx)C_b(x,dz)\), and \(\nu^{(b),c}=\nu^{(b)}C_b\), and
let \(\mathsf U_b:L^2(\nu^{(b)})\to L^2(\nu^{(b),c})\) be the reverse conditional expectation of
\Cref{thm:rg-action-lp-contraction} for this arrow. For every \(h\in L^2_0(\nu)\) and
every path with score \(h\), the pushed path \(\nu_t^{\otimes b}C_b\) is differentiable
in quadratic mean at \(\nu^{(b),c}\) with score
\begin{equation}
\mathscr L_bh:=\mathsf U_b\mathscr I_bh
=\E_{\mathbb J_b}\left[\sum_{i=1}^bh(x_i)\middle| z\right],
\label{eq:rg-score-renormalization-operator}
\end{equation}
so \(\mathscr L_b:L^2_0(\nu)\to L^2_0(\nu^{(b),c})\) is linear with
\(\lVert\mathscr L_b\rVert\leq\sqrt b\). The information budget is exact:
\begin{equation}
b\lVert h\rVert_{L^2(\nu)}^2
-\lVert\mathscr L_bh\rVert_{L^2(\nu^{(b),c})}^2
=\int\operatorname{Var}_{\mathbb J_b}\left(\mathscr I_bh\middle| z\right)
\nu^{(b),c}(dz)
\geq0,
\label{eq:rg-extensive-fisher-defect}
\end{equation}
with equality exactly when \(\mathscr I_bh\) agrees \(\mathbb J_b\)-almost surely with a
measurable function of \(z\). Equivalently, \(b^{-1/2}\mathscr L_b\) contracts
\(L^2_0\) norms. \status{ESTABLISHED}

\noindent\emph{Proof.}
Suppose first that \(h\) is bounded and that the path is
\eqref{eq:rg-quadratic-score-path}. For \(\lvert t\rvert\) small, \(w_t\geq1-\lvert t\rvert\lVert h\rVert_\infty/2>0\),
so \(d\nu_t/d\nu=w_t^2/Z_t=1+th+O(t^2)\) uniformly, the product density is
\(1+t\mathscr I_bh+O(t^2)\) uniformly, and, exactly as in
\eqref{eq:rg-conditional-partition},
\begin{equation}
\frac{d(\nu_t^{\otimes b}C_b)}{d\nu^{(b),c}}(z)
=\E_{\mathbb J_b}\left[
\frac{d\nu_t^{\otimes b}}{d\nu^{(b)}}(x)\middle| z\right]
=1+t\mathscr L_bh(z)+O(t^2)
\label{eq:rg-pushed-density-expansion}
\end{equation}
uniformly. Since \(\sqrt{1+v}=1+v/2+O(v^2)\) uniformly for \(\lvert v\rvert\leq1/2\), the
square root of \eqref{eq:rg-pushed-density-expansion} is
\(1+\tfrac t2\mathscr L_bh+O(t^2)\) uniformly, hence in \(L^2(\nu^{(b),c})\).

For a general path \(\nu_t\) with score \(h\in L^2_0(\nu)\), choose bounded centered
\(h_n\to h\) in \(L^2(\nu)\) and let \(\nu_t^{(n)}\) be their paths
\eqref{eq:rg-quadratic-score-path}. Squared Hellinger distance is the \(f\)-divergence of
\(f(v)=(\sqrt v-1)^2\), so conditional Jensen applied to the disintegration of the
pushed joint law gives
\begin{equation}
\left\lVert
\sqrt{\frac{d(\nu_t^{\otimes b}C_b)}{d\nu^{(b),c}}}
-\sqrt{\frac{d((\nu_t^{(n)})^{\otimes b}C_b)}{d\nu^{(b),c}}}
\right\rVert_{2}
\leq
\left\lVert
\sqrt{\frac{d\nu_t^{\otimes b}}{d\nu^{(b)}}}
-\sqrt{\frac{d(\nu_t^{(n)})^{\otimes b}}{d\nu^{(b)}}}
\right\rVert_{2}.
\label{eq:rg-hellinger-contraction}
\end{equation}
Both input paths have scores \(\mathscr I_bh\) and \(\mathscr I_bh_n\) by
\Cref{prop:rg-score-block-lift}, so the right side divided by \(\lvert t\rvert\) has limit
superior \(\tfrac12\lVert\mathscr I_b(h-h_n)\rVert_2=\tfrac{\sqrt b}2\lVert h-h_n\rVert_2\).
Adding and subtracting the bounded-case expansion and using
\(\lVert\mathscr L_b(h-h_n)\rVert_2\leq\sqrt b\lVert h-h_n\rVert_2\),
\begin{equation}
\limsup_{t\to0}\frac1{\lvert t\rvert}
\left\lVert
\sqrt{\frac{d(\nu_t^{\otimes b}C_b)}{d\nu^{(b),c}}}
-1-\frac t2\mathscr L_bh
\right\rVert_2
\leq\sqrt b\lVert h-h_n\rVert_2 .
\label{eq:rg-pushed-score-approximation}
\end{equation}
The left side does not depend on \(n\), so it is zero. Conditional Jensen gives
\(\lVert \mathsf U_b\rVert\leq1\) and preservation of the mean, so \(\mathscr L_b\) has the stated
type and \(\lVert\mathscr L_b\rVert\leq\lVert\mathscr I_b\rVert\leq\sqrt b\). Finally the
conditional variance identity \eqref{eq:rg-action-total-conditional-variance} applied to
\(\mathscr I_bh\), together with \eqref{eq:rg-score-replication-norm}, is
\eqref{eq:rg-extensive-fisher-defect}, and a nonnegative conditional variance integrates
to zero exactly when it vanishes almost everywhere. \hfill$\square$

\medskip
Both ends of the budget are attained in simple cases. A channel that maps every block
to one point sends every score to zero, so a fixed channel can annihilate the whole
tangent. In the other direction, equality in \eqref{eq:rg-extensive-fisher-defect} at
one parameter is local score sufficiency and not recovery of the experiment, as the
witness following \Cref{thm:cg-fisher-contraction} shows.

\section{Bounded actions and ambient score tangents}
\label{sec:ig-score-charts}

The bounded action calculus of the main manuscript perturbs a law \(\nu\) by bounded
actions \(\varphi\in\mathfrak B=L^\infty(\nu;\R)\) and works on the quotient
\(\overline{\mathfrak B}=\mathfrak B/\R\mathbf 1\) with the oscillation norm
\eqref{eq:rg-action-oscillation-norm}. Its normalized exponential path
\(t\mapsto\widehat\nu^{t\varphi}\) is \eqref{eq:rg-normalized-perturbation-action}. For a
coarse channel with pushed law \(\nu^c\) and reverse conditional expectation \(\mathsf U\),
\(\overline{\mathsf U}\) denotes the descent of \(\mathsf U\) to the quotient.

\propositionheading{Bounded action classes embed isometrically in the Fisher tangent}{prop:rg-action-score-isometry}
For \(\varphi\in\mathfrak B\), the path \(t\mapsto\widehat\nu^{t\varphi}\) is
differentiable in quadratic mean with score
\begin{equation}
\mathscr S_\nu[\varphi]=-\varphi+\nu(\varphi),
\label{eq:rg-action-score-map}
\end{equation}
which depends only on the class \([\varphi]\in\overline{\mathfrak B}\). With the Fisher
norm
\begin{equation}
\lVert[\varphi]\rVert_F
=\inf_{c\in\R}\lVert\varphi-c\rVert_{L^2(\nu)}
=\lVert\varphi-\nu(\varphi)\rVert_{L^2(\nu)},
\label{eq:rg-action-fisher-norm}
\end{equation}
the map \(\mathscr S_\nu:(\overline{\mathfrak B},\lVert\cdot\rVert_F)\to L^2_0(\nu)\) is a
linear isometry onto the bounded centered subspace; this range is dense and, whenever
\(L^2_0(\nu)\) contains an unbounded element, proper, so the Fisher-norm completion of
\(\overline{\mathfrak B}\) is canonically \(L^2_0(\nu)\). The two quotient norms satisfy
\(\lVert[\varphi]\rVert_F\leq\lVert[\varphi]\rVert_{\mathrm{osc}}\). For a coarse channel
the centering square commutes,
\begin{equation}
\mathscr S_{\nu^c}\bigl(\overline{\mathsf U}[\varphi]\bigr)
=\mathsf U\mathscr S_\nu[\varphi],
\label{eq:rg-score-action-square}
\end{equation}
and the replicated block action \((x_1,\dots,x_b)\mapsto\sum_i\varphi(x_i)\) has score
\(\mathscr I_b\mathscr S_\nu[\varphi]\) at \(\nu^{\otimes b}\) and, after the block
channel of \Cref{thm:rg-score-pushforward-defect}, pushed score
\(\mathscr L_b\mathscr S_\nu[\varphi]\). \status{ESTABLISHED}

\noindent\emph{Proof.}
Boundedness gives the uniform expansions \(e^{-t\varphi}=1-t\varphi+O(t^2)\) and
\(\nu(e^{-t\varphi})=1-t\nu(\varphi)+O(t^2)\), so the density ratio of
\(\widehat\nu^{t\varphi}\) is \(1-t[\varphi-\nu(\varphi)]+O(t^2)\) uniformly and its square
root is \(1+\tfrac t2\mathscr S_\nu[\varphi]+O(t^2)\), which is
\eqref{eq:rg-dqm-definition}. Adding a constant to \(\varphi\) leaves
\eqref{eq:rg-action-score-map} unchanged. The \(L^2\) minimizer over constants is the
mean, which gives \eqref{eq:rg-action-fisher-norm} and the isometry, and
\(\inf_c\lVert\varphi-c\rVert_2\leq\inf_c\lVert\varphi-c\rVert_\infty\) gives the norm
order. Bounded functions are dense in \(L^2(\nu)\) and centering is continuous, so the
range is dense; it misses every unbounded element of \(L^2_0(\nu)\), for example
\(h(x)=x\) when \(\nu=\mathcal N(0,1)\). Identity \eqref{eq:rg-score-action-square} is mean
preservation of \(\mathsf U\) together with linearity, and the replication statement is
\Cref{prop:rg-score-block-lift} with \eqref{eq:rg-score-renormalization-operator}.
\hfill$\square$

\medskip
The two norms on \(\overline{\mathfrak B}\) do different work. The oscillation norm
controls the nonlinear map of \Cref{thm:rg-bounded-action-calculus}; the Fisher norm
controls the quadratic-mean tangent. Completing in the Fisher norm introduces
unbounded scores on which the nonlinear bounded action map is not defined, so a
spectral statement proved on \(L^2_0\) is not a statement about the nonlinear bounded
action chart, and conversely. At a Gaussian law the nesting of the three tangent
objects is exact.

\remarkheading{Three tangent objects at a Gaussian law}{rem:ig-three-tangents}
At a nondegenerate Gaussian law \(\nu\) on \(\R^n\), (a) the tangent space of the
Gaussian family, spanned by the scores \eqref{eq:ig-score}, is a subspace of
\(L^2_0(\nu)\) of dimension \(n+n(n+1)/2\); (b) \(L^2_0(\nu)\) is infinite dimensional and,
by \Cref{lem:rg-dqm-realization}, consists exactly of the scores of paths
differentiable in quadratic mean; and (c) the bounded action quotient embeds
isometrically in \(L^2_0(\nu)\) with dense proper range by
\Cref{prop:rg-action-score-isometry}. The tangent space (a) meets the range of (c) only
in \(0\). \status{ESTABLISHED}

\noindent\emph{Proof.}
The score \eqref{eq:ig-score} is a polynomial of degree at most two whose quadratic part
is \(\tfrac12y^\top\Lambda\dot\Sigma\Lambda y\) and whose linear part is
\(\dot\mu^\top\Lambda y-\mu^\top\Lambda\dot\Sigma\Lambda y\). The quadratic part vanishes
only for \(\dot\Sigma=0\), because \(\Lambda\) is invertible and \(\dot\Sigma\) is
symmetric, and then the linear part \(\dot\mu^\top\Lambda y\) vanishes only for
\(\dot\mu=0\); so the dimension is that of the moment chart. The
monomials are linearly independent in \(L^2(\nu)\), so \(L^2_0(\nu)\) is infinite
dimensional. A nonzero element of (a) is a nonconstant polynomial and hence unbounded,
while the range of (c) consists of bounded functions. \hfill$\square$

\medskip
The gap between (b) and (c) is visible along the Hermite directions of the standard
Gaussian.

\propositionheading{Exact two-sided domain of a Hermite exponential action}{prop:ig-hermite-exponential-domain}
Let \(\nu=\mathcal N(0,1)\), let \(\mathrm{He}_k\) be the probabilists' Hermite polynomial
of degree \(k\geq1\), and set
\begin{equation}
Z_k(t)=\int e^{-t\mathrm{He}_k(x)}\nu(dx)\in(0,+\infty].
\label{eq:ig-hermite-normalizer}
\end{equation}
Then \(Z_1(t)<\infty\) for every real \(t\); \(Z_2(t)<\infty\) exactly for
\(t>-\tfrac12\), where \(Z_2(t)=e^t(1+2t)^{-1/2}\); for odd \(k\geq3\),
\(Z_k(t)=+\infty\) for every \(t\neq0\); and for even \(k\geq4\),
\(Z_k(t)<\infty\) exactly for \(t\geq0\). Consequently every \(\mathrm{He}_k\) with
\(k\geq1\) lies in \(L^2_0(\nu)\) and is realized as a two-sided quadratic-mean score,
while for \(k\geq3\) no two-sided exponential action neighborhood of the origin contains
it; the lowest-degree direction whose exponential normalizer diverges on both sides is
\(\mathrm{He}_3(x)=x^3-3x\), with \(\E_\nu\mathrm{He}_3^2=6\). \status{ESTABLISHED}

\noindent\emph{Proof.}
For \(k=1\) the integral is the moment-generating function \(e^{t^2/2}\). For \(k=2\),
\(\mathrm{He}_2(x)=x^2-1\), so \(Z_2(t)=e^t\int e^{-tx^2}\nu(dx)\), which converges
exactly when \(t+\tfrac12>0\), with value \(e^t(1+2t)^{-1/2}\). For \(k\geq3\) write
\(\mathrm{He}_k(x)=x^k+O(x^{k-2})\), so the exponent \(-t\mathrm{He}_k(x)-x^2/2\) of the
integrand against Lebesgue measure has leading term \(-tx^k\). For odd \(k\) and
\(t\neq0\), exactly one tail has \(-tx^k\to+\infty\), and the integral diverges; for
\(\mathrm{He}_3\) and \(t>0\), for instance, \(-t\mathrm{He}_3(-r)-r^2/2=tr^3-r^2/2-3tr\to+\infty\)
as \(r\to+\infty\). For even \(k\) and \(t<0\) both tails diverge; for \(t>0\) the
integrand is eventually bounded by \(e^{-\mathrm c_1x^k}\) with \(\mathrm c_1>0\), and
\(Z_k(0)=1\). Hermite polynomials are square integrable against \(\nu\) and
orthogonal to the constants, the realization is \Cref{lem:rg-dqm-realization}, and
\(\E_\nu\mathrm{He}_3^2=3!\). \hfill$\square$

\medskip
The two score directions of the Gaussian family lie in the benign part of this
classification. At \(n=1\), \(\mu=0\), and \(\Lambda=1\), the score \eqref{eq:ig-score}
reads \(\dot\mu x+\tfrac{\dot\Sigma}2(x^2-1)=\dot\mu\mathrm{He}_1+\tfrac{\dot\Sigma}2\mathrm{He}_2\),
so the mean and variance directions occupy Hermite degrees one and two and both admit
two-sided exponential neighborhoods; the constraint \(t>-\tfrac12\) in the variance
direction is the natural-domain boundary \(\Lambda\succ0\) seen through the exponential
chart. A complete Hermite basis of \(L^2_0(\nu)\) is therefore neither a chart of the
Gaussian family nor a set of directions of the bounded action space, and the spectral
statement below concerns the \(L^2_0\) tangent.

\section{The Gaussian block spectrum}

\theoremheading{Gaussian block spectrum in the normalized Hermite basis}{thm:rg-gaussian-hermite-spectrum}
Let \(\nu=\mathcal N(0,1)\), let \(b\geq2\) be an integer, let \(\xi_1,\dots,\xi_b\) be
independent with law \(\nu\), and let the block channel be the deterministic
normalized sum
\begin{equation}
\bar\xi=b^{-1/2}\sum_{i=1}^b\xi_i,
\label{eq:rg-gaussian-block-statistic}
\end{equation}
whose law is again \(\nu\). That normalization, and no general principle, makes source
and target the same space and hence makes
\begin{equation}
\mathscr L_b=\mathsf U_b\mathscr I_b:L^2_0(\nu)\longrightarrow L^2_0(\nu)
\label{eq:rg-gaussian-score-operator}
\end{equation}
an endomorphism. With \(\widehat{\mathrm{He}}_k=\mathrm{He}_k/\sqrt{k!}\), the family
\(\{\widehat{\mathrm{He}}_k\}_{k\geq1}\) is a complete orthonormal basis of \(L^2_0(\nu)\) and
\begin{equation}
\mathscr L_b\widehat{\mathrm{He}}_k=b^{1-k/2}\widehat{\mathrm{He}}_k,
\qquad k\geq1.
\label{eq:rg-hermite-eigenvalues}
\end{equation}
Consequently \(\mathscr L_b\) is self-adjoint, positive, and diagonal in that basis; it
is Hilbert--Schmidt with \(\sum_{k\geq1}b^{2-k}=b^2/(b-1)\), hence compact; its norm
and spectral radius are \(\sqrt b\); and
\begin{equation}
\sigma(\mathscr L_b)
=\left\{b^{1-k/2}:k\geq1\right\}\cup\{0\}.
\label{eq:rg-hermite-spectrum}
\end{equation}
Every displayed nonzero value is a simple eigenvalue. The value \(0\) is not an
eigenvalue; it is the unique accumulation point of the spectrum and belongs to the
continuous spectrum. \status{ESTABLISHED}

\noindent\emph{Proof.}
Completeness of the normalized Hermite system in \(L^2(\nu)\) is standard, and deleting
\(\widehat{\mathrm{He}}_0=1\) leaves a complete orthonormal basis of the centered subspace. For each
\(i\), the pair \((\xi_i,\bar\xi)\) is jointly centered Gaussian with unit variances and
covariance \(b^{-1/2}\), so \(\xi_i\) given \(\bar\xi=z\) has law
\(\mathcal N(z/\sqrt b,1-1/b)\). Taking the conditional expectation of the generating
function \(\exp(t\xi_i-t^2/2)\) and completing the square,
\begin{equation}
\E\left[e^{t\xi_i-t^2/2}\middle|\bar\xi=z\right]
=\exp\left(\frac{tz}{\sqrt b}+\frac{t^2}2\left(1-\frac1b\right)-\frac{t^2}2\right)
=\exp\left(sz-\frac{s^2}2\right),
\qquad s=\frac t{\sqrt b}.
\label{eq:rg-mehler-generating}
\end{equation}
The generating identity \(\exp(sz-s^2/2)=\sum_{k\geq0}\mathrm{He}_k(z)s^k/k!\) converges in
\(L^2(\nu)\) locally uniformly in \(s\), and conditional expectation is \(L^2\)
continuous, so the two power series in \(t\) may be compared coefficient by
coefficient. This gives the Mehler regression identity
\(\E[\mathrm{He}_k(\xi_i)\mid\bar\xi]=b^{-k/2}\mathrm{He}_k(\bar\xi)\), and summing over the
\(b\) replicated coordinates and dividing by \(\sqrt{k!}\) proves
\eqref{eq:rg-hermite-eigenvalues}, uniformly in \(b\).

Since \(\mathscr L_b\) is bounded and acts diagonally on a complete orthonormal basis, it
is that diagonal operator, hence self-adjoint and positive with eigenvalues
\(b^{1-k/2}\downarrow0\). Square summability of the eigenvalues is the stated geometric
series, so \(\mathscr L_b\) is Hilbert--Schmidt and compact, and its norm is the largest
eigenvalue \(b^{1/2}\). The spectrum of a diagonal operator is the closure of its
eigenvalue set, which is \eqref{eq:rg-hermite-spectrum}, and the eigenvalues are pairwise
distinct, so each is simple. Every eigenvalue is strictly positive, so
\(\ker\mathscr L_b=\{0\}\). The range contains every finite Hermite expansion and is
dense, but it is proper: \(\sum_{k\geq1}b^{1-k/2}\widehat{\mathrm{He}}_k\) lies in \(L^2_0(\nu)\), while any
preimage would have every Hermite coefficient equal to one, which is not square
summable. An injective operator with dense, nonclosed range has \(0\) in its continuous
spectrum. \hfill$\square$

\medskip
The case \(b=2\) of \eqref{eq:rg-hermite-eigenvalues} is the linearization of the
normalized-sum map on densities at the standard Gaussian computed by
\citet[Sec.~2, Eq.~(2.14)]{JonaLasinio2001}, who notes there that completing a limit
theorem requires control of the nonlinear terms, which he does not pursue. The
theorem above is a statement about the derivative at the Gaussian law, a linear
operator on the \(L^2_0\) tangent; the eigenvalue \(b^{1/2}\) of the degree-one direction
exceeds one only because the replication factor of
\eqref{eq:rg-score-replication-norm} is extensive while the norm is held fixed, since
the one-arrow operator \(\mathsf U_b\) is a contraction.
